\begin{abstract}
    \noindent
    Replay-based continual learning is usually judged by the accuracy it reaches \emph{after} each task. Continual evaluation exposes a failure this view hides: the \emph{stability gap}, a sharp but transient drop in past-task accuracy at the start of every new task. The gap persists even when the optimiser descends the exact joint loss over all tasks, so its cause lies in the optimisation trajectory rather than the objective. This thesis explains the gap on two levels. Learning both tasks has a fixed price: the loss on past tasks must rise to its level at the joint optimum. An envelope-theorem argument shows that a reference path exists along which this rise is monotone, so everything the past-task loss does above that level is avoidable in principle. Steepest descent after an abrupt task switch leaves that path: it bows out of the old task's basin and returns from outside, a deterministic arc that survives exact gradients (level one, the trajectory problem). Real optimisers also overshoot whatever path they follow, pushed by an imbalanced first step and by gradient noise from a finite buffer (level two, the overshoot problem). Written in a unified update rule, every replay method counters these problems with the same tool, a \emph{gate} on the new-task gradient, while the replay gradient restores at full strength. We build and compare the two pure gate designs. The \emph{$\lambda$-curriculum} is a scalar feedback gate: it ramps the new-task loss weight up from zero, so the optimiser tracks the moving valley floor of the reference path. \emph{Asymmetric preconditioned experience replay} is a directional feedforward gate: it filters the new-task gradient through past-task curvature, passing it where the past task is flat and braking it where the past task is steep. On rotated MNIST the two levels separate cleanly: removing the overshoot drivers leaves the arc, placing the optimiser on the reference path leaves the overshoot, and doing both removes the gap while the fixed price remains. The curriculum needs momentum: its acceleration repays the plasticity the schedule borrows, and its averaging supplies the variance reduction that closes the residual gap. The directional gate lowers the gap monotonically as its filter strengthens, at no plasticity cost but higher compute; at its strongest setting it fails late and abruptly, consistent with blind spots of its sampled curvature metric, and momentum re-inflates its per-step filter. Which gate wins is therefore set by the momentum regime: without momentum the directional gate is better on both stability and accuracy, but with momentum, which realistic training requires, the scalar gate takes over. On CIFAR-10 with a ResNet-18 the curriculum cuts gap depth several-fold below vanilla experience replay at matched final accuracy and negligible compute cost, under two normalisation schemes and on corruptions it was never tuned on.
\end{abstract}
